信息论进入机器学习之后，往往被压缩成三句口号：信息增益用来选特征，交叉熵用来训练网络，信息瓶颈用来解释表示学习。三句都不算错，可它们各自都省掉了成立所需的前提，而省掉的那部分恰恰是判断力所在。

先看一个数值反例。决策树的父节点有八个样本，正负各四，熵是 \(1\) 比特。某个划分把正负完美分开，信息增益 \(1\) 比特；另一个划分得到 \((3+,1-)\) 与 \((1+,3-)\)，两边的熵都是 \(h_2(0.75)\approx0.811\)，信息增益只有约 \(0.189\) 比特。前者显然更好。可如果把样本编号当特征用，每个叶子只装一个样本，信息增益也是 \(1\) 比特——同一个满分，一个记录规律，一个记录背诵。互信息量的是关联，它既不筛因果，也不防过拟合。

这一章要立住四条判断，它们是这本书``面向 AI 时代''这个定位最直接的落点。

其一，分类训练时最小化交叉熵，等价于最小化模型与真实条件分布之间的 KL 散度，因为真实分布的熵是一个与参数无关的常数。日常调的那条损失曲线，压不下去的水平段就是数据自身的熵。

其二，最小描述长度把模型复杂度与数据拟合放到同一把尺子——比特——上量。这样一来，正则化项不再是外加的惩罚，它就是``说明这个模型''所需的长度，而它与 Bayesian 的对应关系相当精确：负对数先验就是模型的编码长度。

其三，信息瓶颈给出了``什么叫好表示''的一个可以写下来的定义：压缩输入、保留标签，目标是 \(\min I(X;T)-\beta I(T;Y)\)。同时必须说清它的边界——确定性网络里 \(I(X;T)\) 可能是无穷或常数，读数依赖离散化与估计器，相关的实验结论至今有争议。

其四，Fano 不等式把条件熵翻译成错误率的硬下界。特征里的信息不足以覆盖标签的不确定性时，任何算法的错误率都有一个够不着的地板。这条判断在实践中最值钱，因为它区分的是``模型还不够好''和``信息本来就不够''，而这两种诊断指向完全相反的行动。

\begin{center}
\begin{tikzpicture}[x=1.32cm,y=1cm,font=\small]
  \draw[->,black!55] (0,0) -- (10.1,0) node[right,font=\footnotesize] {建模流程};
  \foreach \x in {1.2,4.0,6.8,9.2} \draw[black!45] (\x,-0.12) -- (\x,0.12);
  \fill[blue!60] (1.2,0) circle (2.2pt);
  \fill[blue!60] (4.0,0) circle (2.2pt);
  \fill[blue!60] (6.8,0) circle (2.2pt);
  \fill[red!70] (9.2,0) circle (2.2pt);
  \node[font=\footnotesize,align=center,black!80] at (1.2,0.78) {该收哪份数据};
  \node[font=\footnotesize,align=center,black!80] at (4.0,0.78) {模型该多复杂};
  \node[font=\footnotesize,align=center,black!80] at (6.8,0.78) {表示该留什么};
  \node[font=\footnotesize,align=center,black!80] at (9.2,0.78) {还剩多少余地};
  \node[font=\footnotesize,align=center,black!60] at (1.2,-0.62) {\(I(Y;X)\)};
  \node[font=\footnotesize,align=center,black!60] at (4.0,-0.62) {\(L(M)+L(D\given M)\)};
  \node[font=\footnotesize,align=center,black!60] at (6.8,-0.62) {\(I(X;T)-\beta I(T;Y)\)};
  \node[font=\footnotesize,align=center,black!60] at (9.2,-0.62) {\(P_e\) 的下界};
\end{tikzpicture}

\emph{四篇分别站在建模流程的四个位置上，用的却是同一个单位。前三个位置回答``该怎么选''，最后一个位置回答``选完之后还有没有希望''——它是唯一一条能告诉你该停手的判断。}
\end{center}

\subsection{怎么读这一章}

第一次读先抓问题和反例：每一篇都配了一个能拆开算的最小例子，跟着算一遍比记住公式管用。遇到互信息估计、连续确定性网络或者因果解释，再回头核对随机变量和概率模型究竟是怎样定义的——这一章里几乎所有的争议，追到底都是定义没对齐。

需要的前置知识集中在两处：互信息与数据处理不等式来自\hyperref[ch:055]{``联合信息：一个变量能告诉我们另一个变量多少''}，交叉熵与 KL 的分解来自\hyperref[ch:056]{``分布之间的差异：从错误编码到散度''}。上一章是\hyperref[ch:058]{``有噪信道编码：噪声中为什么仍能可靠通信''}，那里的容量概念在本章会以``实验值多少比特''的身份重新出现一次。

\subsection{本章篇目}

\begin{itemize}
\tightlist
\item \hyperref[sec:07-Information-Theory/06-Information-and-Learning/01]{``信息增益：从提问到实验设计''}：给``下一次观察值不值得''定价，并说明为什么经验信息增益总是往高了报；
\item \hyperref[sec:07-Information-Theory/06-Information-and-Learning/02]{``最小描述长度与 Bayesian 学习''}：把模型复杂度换算成比特，于是正则化、MAP 与 AIC/BIC 落在同一张账本上；
\item \hyperref[sec:07-Information-Theory/06-Information-and-Learning/03]{``信息瓶颈与任务相关表示''}：把``好表示''写成两个互信息的取舍，同时交代它在确定性网络上的定义困难；
\item \hyperref[sec:07-Information-Theory/06-Information-and-Learning/04]{``分类、神经网络与信息论边界''}：交叉熵的地板、困惑度的换算，以及 Fano 界怎样区分模型不够好与信息不够。
\end{itemize}

四篇按概念依赖排列，但从具体困惑进入同样可行：正卡在特征选择上就从第一篇读起，正在为模型选择准则头疼就直接翻第二篇，而如果手头的指标已经很久不动了，第四篇是最该先读的一篇。
